\usepackage{xspace}
\usepackage{amsmath}
\usepackage{amsfonts}
\usepackage{amssymb}
\usepackage{bm}
\usepackage{upgreek}
\usepackage{graphicx}
\usepackage{ifthen}


\newcommand{\latex}{\LaTeX\xspace}


%%%%%%%%%%%%%%%%
%              %
%  references  %
%              %
%%%%%%%%%%%%%%%%

\newcommand{\reffig}[1]{Figure~\ref{fig:#1}}
\newcommand{\refsubfig}[1]{Figure~\ref{subfig:#1}}


%%%%%%%%%%
%        %
%  sets  %
%        %
%%%%%%%%%%


\newcommand{\setN}{\ensuremath{\mathbb{N}}}  % natural numbers
\newcommand{\setZ}{\ensuremath{\mathbb{Z}}}  % integers
\newcommand{\setQ}{\ensuremath{\mathbb{Q}}}  % rational numbers
\newcommand{\setR}{\ensuremath{\mathbb{R}}}  % reals
\newcommand{\setC}{\ensuremath{\mathbb{C}}}  % complex numbers



%%%%%%%%%%%%%%%%%%%%%%%%%%%%%
%                           %
%  formatting of variables  %
%                           %
%%%%%%%%%%%%%%%%%%%%%%%%%%%%%

% variable:
\newcommand{\var}[1]{\ensuremath{#1}}

% complex variable:
\newcommand{\cvar}[1]{\ensuremath{\underline{#1}\,}}

% vectors
\renewcommand{\vec}[1]{\ensuremath{\bm{#1}}}
\newcommand{\vecTr}[1]{\ensuremath{\bm{#1}\mTr}}
\newcommand{\vecHerm}[1]{\ensuremath{\bm{#1}}\mHerm}

\newcommand{\cvec}[1]{\vec{\cvar{#1}}}
\newcommand{\cvecTr}[1]{\vecTr{\cvar{#1}}}
\newcommand{\cvecHerm}[1]{\vecHerm{\cvar{#1}}}


% matrices
\newcommand{\mat}[1]{\vec{#1}}
\newcommand{\matTr}[1]{\vecTr{#1}}
\newcommand{\matHerm}[1]{\vecHerm{#1}}

\newcommand{\cmat}[1]{\cvec{#1}}
\newcommand{\cmatTr}[1]{\cvecTr{#1}}
\newcommand{\cmatHerm}[1]{\cvecHerm{#1}}



%%%%%%%%%%%%%%%%%%%%%%%%%%%%%
%                           %
%  formatting of functions  %
%                           %
%%%%%%%%%%%%%%%%%%%%%%%%%%%%%


\newcommand{\func}[1]{\ensuremath{\mathrm{#1}}}                          % function that returns a real-valued scalar
\newcommand{\cfunc}[1]{\ensuremath{\underline{\mathrm{#1}}\,}}           % function that returns a complex-valued scalar
\newcommand{\vecfunc}[1]{\ensuremath{\bm{\mathrm{#1}}}}                  % function that returns a real-valued vector
\newcommand{\cvecfunc}[1]{\ensuremath{\underline{\bm{\mathrm{#1}}}\,}}   % function that returns a complex-valued vector
\newcommand{\matfunc}[1]{\vecfunc{#1}}                                   % function that returns a real-valued matrix
\newcommand{\cmatfunc}[1]{\cvecfunc{#1}}                                 % function that returns a complex-valued scalar



%%%%%%%%%%%%%%%%%%%%%%%%%%%%%%%%%%%%%%
%                                    %
%  standard functions and operators  %
%                                    %
%%%%%%%%%%%%%%%%%%%%%%%%%%%%%%%%%%%%%%


% operators on complex numbers
\newcommand{\cConj}{\ensuremath{^\ast}}                                             % complex conjugate
\newcommand{\cReal}[1]{\ensuremath{\mathcal{R}\left\lbrace #1 \right\rbrace}}      % real part
\newcommand{\cImag}[1]{\ensuremath{\mathcal{I}\left\lbrace #1 \right\rbrace}}      % imaginary part


% unary matrix/vector operators (suffixes):
\newcommand{\mTr}{\ensuremath{^\mathrm{T}}}           % transposition
\newcommand{\mHerm}{\ensuremath{^\mathrm{H}}}         % Hermitian transpose (complex conjugate transpose)

\newcommand{\mInv}{\ensuremath{^{\text{-}1}}}                % matrix inverse
\newcommand{\mInvTr}{\ensuremath{^{-\mathrm{T}}}}     % transpose of matrix inverse
\newcommand{\mInvHerm}{\ensuremath{^{-\mathrm{H}}}}   % hermitian transpose of matrix inverse
\newcommand{\mPinv}{\ensuremath{^\text{\dag}}}            % Moore-Penrose pseudo-inverse

\newcommand{\matInv}[1]{\ensuremath{\mat{#1}\mInv}}
\newcommand{\cmatInv}[1]{\ensuremath{\cmat{#1}\mInv}}
\newcommand{\matPinv}[1]{\ensuremath{\mat{#1}\mPinv}}
\newcommand{\cmatPinv}[1]{\ensuremath{\cmat{#1}\mPinv}}


% binary matrix/vector operators:
\newcommand{\mKron}{\ensuremath{\otimes}}             % Kronecker product
\newcommand{\mKhatri}{\ensuremath{\diamond}}          % Khatri-Rao product
\newcommand{\mElemProd}{\ensuremath{\odot}}           % hadamard (element-wise) product


% other operators
\newcommand{\mNorm}[2]{\ensuremath{\left\| #1 \right\|_{#2}}}         % general norm
\newcommand{\mFrobNorm}[1]{\mNorm{#1}{\text{F}}}                      % Frobenius norm
\newcommand{\mDiag}[1]{\ensuremath{\mathrm{diag}\left(#1\right)}}     % for matrices: returns main diagonal as vector; for vecor: returns diagonal matrix
\newcommand{\mTrace}[1]{\ensuremath{\mathrm{trace}\left(#1\right)}}   % sum of diagonal elements
\newcommand{\mDet}[1]{\ensuremath{\mathrm{det}\left(#1\right)}}       % determinant



%%%%%%%%%%%%%%%%%
%               %
%  convolution  %
%               %
%%%%%%%%%%%%%%%%%

\makeatletter
\newcommand{\convolve}{\ensuremath{\ast}}
\newcommand{\circConvolve}{\ensuremath{\mathbin{\mathpalette\make@circled\ast}}}
\newcommand{\make@circled}[2]{%
  \ooalign{$\m@th#1\smallbigcirc{#1}$\cr\hidewidth$\m@th#1#2$\hidewidth\cr}%
}
\newcommand{\smallbigcirc}[1]{%
  \vcenter{\hbox{\scalebox{0.77778}{$\m@th#1\bigcirc$}}}%
}
\makeatother


%%%%%%%%%%%%%%%%%%%%%%%%
%                      %
%  coordinate systems  %
%                      %
%%%%%%%%%%%%%%%%%%%%%%%%

% azimuth angle:
\newcommand{\azi}{\ensuremath{\varphi}}
\newcommand{\upazi}{\ensuremath{\upvarphi}}  % non-italic

% elevation angle:
\newcommand{\ele}{\ensuremath{\vartheta}}
\newcommand{\upele}{\ensuremath{\upvartheta}}  % non-italic

% helper command for command \coele and  \upcoele (see below)
\newcommand{\hslashslash}{%
  \raisebox{.9ex}{%
    \scalebox{.7}{%
      \rotatebox[origin=c]{18}{$-$}%
    }%
  }%
}

% co-elevation:
%   - Ranges from positive y-axis (0°/0 rad) to negative y-axis (180°/? rad).
%   - Formatted using a \vartheta with a slash.
\DeclareRobustCommand{\coele}{%
  {% keep these braces!
    \ensuremath{%                            
    \vphantom{\vartheta}%
    \ooalign{\kern-.07em\smash{\lower.35em\hbox{\hslashslash}}\hidewidth\cr$\vartheta$\cr}%
    \kern.05em
    }%
  }%
}

% non-italic co-elevation
\newcommand{\upcoele}{%
  {% keep these braces!
    \text{% needed to scale correctly in subscripts
      \ensuremath{%
        \vphantom{\upvartheta}%
        \ooalign{\kern-.07em\smash{\lower.35em\hbox{\hslashslash}}\hidewidth\cr$\upvartheta$\cr}%
        \kern.05em%
      }%
    }%
  }%
}



%%%%%%%%%%%%%%%%%
%               %
%  derivatives  %
%               %
%%%%%%%%%%%%%%%%%

% differential 'd':
\newcommand{\diff}{\mathrm{d}}

% make \partial symbol upright:
\let\oldpartial\partial
\DeclareRobustCommand{\uppartial}{\text{\rotatebox[origin=t]{20}{\scalebox{0.95}[1]{$\oldpartial$}}}\hspace{-1pt}}
\renewcommand{\partial}{\uppartial}

\newcommand{\pDiff}{\uppartial} % alias (to be consistent with \diff)


\newcommand{\deriv}[3][]{%
  \ifthenelse{\equal {#1}{}}{% no order given:
    \frac{\diff #2}{\diff #3}%
  }{ % order given:
    \frac{\diff^{#1} #2}{\diff #3^{#1}}%
  }%
}

\newcommand{\pDeriv}[3][]{%
  \ifthenelse{\equal {#1}{}}{% no order given:
    \frac{\pDiff #2}{\partial #3}%
  }{% order given:
    \frac{\pDiff^{#1} #2}{\partial #3^{#1}}%
  }%
}



%%%%%%%%%%%%%%%%%%%%%%%%%%%%%%%%
%                              %
%  general transform commands  %
%                              %
%%%%%%%%%%%%%%%%%%%%%%%%%%%%%%%%

% based on package trfsgns (https://ctan.org/pkg/trfsigns):

% transform:
%
% #1: argument for the resulting function or series
% #2: letter or name of the transform
% #3: expression that is being transformed
% #4: subscript text (e.g., for giving the variable over which the transform in performed)
\newcommand{\mathTransform}[4]{%
  \ensuremath{%
      \operatorname*{#2}%
      \ifthenelse{\equal {#4}{}}{% no argument for subscript text        
      }{% argument for subscript text given
        _{#4} %
      }%
      \kern-2.0pt%
      \left\lbrace #3 \right\rbrace%
      \ifthenelse{\equal {#1}{}}{% no argument for resulting function given        
      }{% argument for resulting function given
        \kern-2.0pt #1 %
      }%
  }%
}


% continues transform
%
% #1: argument for the resulting function (example: F{g(t)}(f) <-- 'f' is the argument)
%     --> If not empty, the argument will be enclosed in round brackets.
% #2: letter or name of the transform
% #3: expression that is being transformed
\newcommand{\mathContTransform}[4]{%
  \mathTransform%
  { \if\relax\detokenize{#1}\relax\else\ensuremath{\left( #1 \right)}\fi }% enclose in round brackets if optional argument is given
  {#2}{#3}{#4}%
}

% discrete transform
%
% #1: argument for the resulting series (example: F{g[n]}[k] <-- 'k' is the argument)
%     --> If not empty, the argument will be enclosed in square brackets.
% #2: letter or name of the transform
% #3: expression that is being transformed
\newcommand{\mathDiscTransform}[4]{%
  \mathTransform%
  { \if\relax\detokenize{#1}\relax\else\ensuremath{\left[ #1 \right]}\fi }% enclose in square brackets if optional argument is given
  {#2}{#3}{#4}%
}

% inverse continues transform
%
% #1: argument for the resulting function (example: F^-1{g(f)}(t) <-- 't' is the argument)
%     --> If given, the argument will be enclosed in round brackets.
% #2: letter or name of the transform
% #3: expression that is being inverse-transformed
\newcommand{\mathContInvTransform}[4]{\mathContTransform{#1}{#2\mInv}{#3}{#4}}

% discrete inverse transform
%
% #1: argument for the resulting series (example: F^-1{g[n]}[k] <-- 'k' is the argument)
%     --> If given, the argument will be enclosed in square brackets.
% #2: letter or name of the transform
% #3: expression that is being inverse-transformed
\newcommand{\mathDiscInvTransform}[4]{\mathDiscTransform{#1}{#2\mInv}{#3}{#4}}


% Transform relation.
% --> Has a a hollow circle at the left and a solid cicle on the right, connected by a horizontal line.
%
% #1 (optional): Symbol for the transformation (written above the horizontal line)
\makeatletter
\newcommand{\transformRelation}[1][]{%
  \mathop{%
    \mbox{%
      \setlength{\unitlength}{0.1em}%
      \begin{picture}(37,10)%
        \put(4,3){\circle{4}}%
        \put(6,3){\line(1,0){26}}%
        \put(33,3){\circle*{4}}%
        \put(19,7){\makebox(0,0){#1}}%
      \end{picture}%
    }%
  }\slimits@%
}
\makeatother


% Reversed transform relation.
% --> Like \transformRelation but with the solid circle at the left and the hollow circle at the right
%
% #1: Symbol for the transformation (written above the horizontal line)
\makeatletter
\newcommand{\revTransformRelation}[1][]{%
  \mathop{%
    \mbox{%
      \setlength{\unitlength}{0.1em}%
      \begin{picture}(37,10)%
        \put(4,3){\circle*{4}}%
        \put(6,3){\line(1,0){25}}%
        \put(33,3){\circle{4}}%
        \put(19,7){\makebox(0,0){#1}}%
      \end{picture}%
    }%
  }\slimits@%
}
\makeatother

% Create a continuous transform
%
% This will create commands \<name>Transform, \<name>InvTransform,
% \<name>Relation, and \<name>RevRelation, where <name> is the name
% of the transform.
%
% #1: name of the transform ('fourier', 'dft', 'laplace', ...)
% #2: Symbol/text used for the transform
\newcommand{\createContinuousTransform}[2]{%
%
  % create command \<#1>Transformation
  \expandafter\NewDocumentCommand\csname#1Transform\endcsname{E{_}{\relax}O{}m}{%
    \mathContTransform{##2}{#2}{##3}{##1}%
  }%
%
  % create command \<#1>InvTransformation
  \expandafter\NewDocumentCommand\csname#1InvTransform\endcsname{E{_}{\relax}O{}m}{%
    \mathContInvTransform{##2}{#2}{##3}{##1}%
  }%
%
  % create command \<#1>Relation
  \expandafter\newcommand\csname#1Relation\endcsname{%
    \mathop{\transformRelation[#2]{}}
  }%
%
  % create command \<#1>RevRelation
  \expandafter\newcommand\csname#1RevRelation\endcsname{%
    \mathop{\revTransformRelation[#2]{}}
  }%
}

% Create a discrete transform
%
% This will create commands \<name>Transform, \<name>InvTransform,
% \<name>Relation, and \<name>RevRelation, where <name> is the name
% of the transform.
%
% #1: name of the transform ('fourier', 'dft', 'laplace', ...)
% #2: Symbol/text used for the transform
\newcommand{\createDiscreteTransform}[2]{%
%
  % create command \<#1>Transformation
  % \expandafter\newcommand\csname#1Transform\endcsname[2][]{%
  \expandafter\NewDocumentCommand\csname#1Transform\endcsname{E{_}{\relax}O{}m}{%
    \mathDiscTransform{##2}{#2}{##3}{##1}%
  }%
%
  % create command \<#1>InvTransformation
  \expandafter\NewDocumentCommand\csname#1InvTransform\endcsname{E{_}{\relax}O{}m}{%
    \mathDiscInvTransform{##2}{#2}{##3}{##1}%
  }%
%
  % create command \<#1>Relation
  \expandafter\newcommand\csname#1Relation\endcsname{%
    \mathop{\transformRelation[#2]{}}
  }%
%
  % create command \<#1>RevRelation
  \expandafter\newcommand\csname#1RevRelation\endcsname{%
    \mathop{\revTransformRelation[#2]{}}
  }%
}


%%%%%%%%%%%%%%%%%%%%%%%%%
%                       %
%  specific transforms  %
%                       %
%%%%%%%%%%%%%%%%%%%%%%%%%


%% Fourier transform %%

% letter for the Fourier transform:
\newcommand{\fourierLetter}{\ensuremath{\mathcal{F}}}

% Create command:
%  \fourierTransform, \fourierInvTransform, \fourierRelation, \fourierRevRelation
\createContinuousTransform{fourier}{\fourierLetter}


%% Discrete Fourier transform %%

% letter for the Fourier transform:
\newcommand{\dftLetters}{\ensuremath{\mathcal{DFT}}}

% Create command:
%  \dftTransform, \dftInvTransform, \dftRelation, \dftRevRelation
\createDiscreteTransform{dft}{\dftLetters}


%% Laplace transform %%

% letter for the Laplace transform:
\newcommand{\laplaceLetter}{\ensuremath{\mathcal{L}}}

% Create command:
%  \laplaceTransform, \laplaceInvTransform, \laplaceRelation, \laplaceRevRelation
\createContinuousTransform{laplace}{\laplaceLetter}



%%%%%%%%%%%%%%%%%
%               %
%  stochastics  %
%               %
%%%%%%%%%%%%%%%%%

\NewDocumentCommand\expValue{e{_}m}{%
    \operatorname*{\mathbb{E}}\IfNoValueF{#1}{_{#1}} \left\lbrace #2 \right\rbrace%
}

\NewDocumentCommand\variance{e{_}m}{%
    \operatorname*{\mathrm{var}}\IfNoValueF{#1}{_{#1}} \left\lbrace #2 \right\rbrace%
}

\NewDocumentCommand\covariance{e{_}m}{%
    \operatorname*{\mathrm{cov}}\IfNoValueF{#1}{_{#1}} \left\lbrace #2 \right\rbrace%
}

\NewDocumentCommand\correlation{e{_}m}{%
    \operatorname*{\mathrm{corr}}\IfNoValueF{#1}{_{#1}} \left\lbrace #2 \right\rbrace%
}

\newcommand{\mCorr}[1]{\ensuremath{\mat{R}_{#1}}}
\newcommand{\mcCorr}[1]{\ensuremath{\cmat{R}_{#1}}}
\newcommand{\mCov}[1]{\ensuremath{\mat{C}_{#1}}}
\newcommand{\mcCov}[1]{\ensuremath{\cmat{C}_{#1}}}